\documentclass{subfiles}
\begin{document}
    We have seen how trees and lists can be used to implent dictionaries.
        Now we discuss another way to implent dictionaries when the set of possible keys \(K\)
        is much greater than the number of key that can be stored.
        In other words, when \(\abs{K} \gg \abs{T}\), 
        where \(T\) is the structure used to store the keys.

    In this section, we first recall the idea behind hashing, 
        we then introduce the notion of universal hash function,
        and show how, under some reasonable assumptions,
        we can achive a \(\ofO{1}\) time for any of the dictionary operations.

    \subsection{What is hashing?}
    \subfile{../sub/3.1 - The idea behind hashing}

    \subsection{Universal hashing}
    \subfile{../sub/3.2 - Universal hashing}

    \subsection{Perfect hashing}
    \subfile{../sub/3.3 - Perfect hashing}
    
    \subsection{Bloom's filter}
    \subfile{../sub/3.4 - Bloom filter}
\end{document}
